\majorheading{MỞ ĐẦU}

Các trò chơi nhập vai có thế giới truyện phức tạp thường khiến người chơi phải tra cứu nhiều trang wiki để hiểu một nhân vật hoặc một sự kiện. Với Elden Ring, phần lớn cốt truyện được kể gián tiếp qua mô tả vật phẩm, lời thoại và địa điểm, nên một câu hỏi đơn giản cũng có thể cần ghép thông tin từ nhiều trang. Các mô hình ngôn ngữ lớn trả lời trôi chảy nhưng có thể đưa ra chi tiết không có trong nguồn. Vì vậy, đề tài kết hợp truy xuất thông tin với mô hình ngôn ngữ để câu trả lời bám vào tài liệu cụ thể và có trích dẫn.


Đề tài xây dựng Index Beast, một hệ thống hỏi đáp về lore Elden Ring bản gốc (base game) bằng tiếng Việt và tiếng Anh, dùng chung một phần phụ trợ cho website và bot Discord. Hệ thống dựa trên kiến trúc sinh văn bản có truy xuất (Retrieval-Augmented Generation, RAG): câu hỏi được dùng để tìm các đoạn văn liên quan trong kho tri thức, sau đó mô hình ngôn ngữ viết câu trả lời dựa trên các đoạn đó. Việc tìm kiếm kết hợp biểu diễn vector ngữ nghĩa với tìm kiếm từ khóa, rồi hợp nhất và xếp hạng lại kết quả.


Mục tiêu của báo cáo là trình bày cơ sở lý thuyết của các thành phần, phát biểu bài toán, mô tả dữ liệu và thiết kế hệ thống, đồng thời ghi lại kết quả của giai đoạn đầu (Sprint 1 trong kế hoạch sáu sprint). Giai đoạn này mới hoàn thành nền tảng kỹ thuật, quy trình thu thập dữ liệu, bộ câu hỏi đánh giá và một thí nghiệm cơ sở trên tập con dữ liệu. Báo cáo không khẳng định chất lượng của hệ thống hoàn chỉnh; các số liệu được nêu cùng điều kiện và giới hạn của từng lần đo.


Báo cáo gồm ba chương. Chương 1 trình bày cơ sở lý thuyết về xử lý ngôn ngữ tự nhiên, truy xuất thông tin, mô hình ngôn ngữ lớn và đánh giá. Chương 2 trình bày bài toán, dữ liệu, phương pháp và thiết kế hệ thống. Chương 3 trình bày thử nghiệm, kết quả, phân tích và hướng phát triển.
